% Input this file after defining theorem, lemma, definition, and proof environments.
% Main document should supply amsmath, amssymb, amsthm, hyperref.
\section{The restricted-witness theorem}
\label{sec:restricted-witness}

This section uses raw iterator terms, with ordinary capture-avoiding
substitution.  No computation is built into the syntax.  Write
\[
 C(v,a,b):=I(v,a;u.b),\qquad 1:=S0,
\]
where the bound variable $u$ does not occur in $b$.  Thus
$C(0,a,b)=a$ and $C(Sv,a,b)=b$ are derivable computation equations.
Formulas are equations, $Q_y A$, and $D_{x,y}A$, with their unchanged
ordinary readings
\[
 Q_yA=A(0)\to\forall y\,A(y),\qquad
 D_{x,y}A=\exists x\bigl(A(x,0)\to\forall y\,A(x,y)\bigr).
\]
The distinction between a variable and the numerical constant $0$ is
syntactic throughout.

\subsection{Seven primitive rule patterns}

Let $\mathcal R$ have Hyp, the two iterator computation equations, the
one-term crossed transport rule
\[
 \frac{\Gamma\vdash s=t\qquad\Gamma\vdash h(s)=r}
      {\Gamma\vdash r=h(t)},
\]
and the usual two eliminations
\[
 \frac{\Gamma\vdash Q_yA(y)\qquad\Gamma\vdash A(0)}
      {\Gamma\vdash A(t)},
 \qquad
 \frac{\Gamma\vdash D_{x,y}B(x,y)\qquad
       \Gamma,Q_yB(b,y)\vdash F}
      {\Gamma\vdash F}.
\]
In the second rule $b$ is fresh in the conclusion and the undischarged
context.  Its one introduction rule is
\[
 \frac{\Gamma,B(t,\eta)\vdash B(t,Sz)}
      {\Gamma\vdash D_{x,y}B(x,y)},
 \qquad
 \eta=\begin{cases}
 z,&B\text{ is a literal equation and }t\text{ is the literal }0,\\
 0,&\text{otherwise},
 \end{cases}
 \tag{RDI}
\]
where $z$ is fresh and the supplied witness $t$ must be either the literal
$0$ or a variable.  In particular, every variable witness selects cases.
There is no separate rule for introducing $Q$, and no general term
substitution rule on proofs.

\begin{lemma}[Structural rules available before completeness]
Context weakening, injective renaming of numerical variables, and cut are
admissible in $\mathcal R$.
\end{lemma}
\begin{proof}
Each assertion follows by induction on the given derivation.  A numerical
renaming sends a variable witness to a variable witness and fixes the
literal $0$; it neither changes the root constructor of a formula nor the
introduction mode.  Renamings are extended under each binder by choosing a
fresh variable.  Context weakening does not change a witness.  In the cut
induction the proof replacing a discharged hypothesis is weakened and
injectively renamed under eigenvariables; all other rule data are
unchanged.  This argument does not assert closure under substitution of an
arbitrary term for a numerical variable.  That property will follow only
after unrestricted introduction has been recovered.
\end{proof}

Reflexivity and the usual equational rules are already derived.  For
example, transport $I(0,a;u.u)=a$ along itself, with the identity template,
to obtain $a=a$.  Applying crossed transport with major premise $s=s$,
minor premise $s=t$, and the identity template gives $t=s$.
Transitivity and congruence follow
from crossed transport and reflexivity.  These constructions need no
logical introduction.

\subsection{Unrestricted case introduction for a closure}

\begin{lemma}[Closure introduction]
Both of the following rules are admissible:
\[
 \frac{\Gamma,A(0)\vdash A(Sz)}{\Gamma\vdash Q_yA(y)},
 \qquad
 \frac{\Gamma,E(z)\vdash E(Sz)}{\Gamma\vdash Q_yE(y)},
 \tag{QC, QP}
\]
where $z$ is fresh and $E$ is a literal equation.  The first rule permits
an arbitrary native formula $A$, including an equation.
\end{lemma}
\begin{proof}
For QP use RDI at witness $0$ on $D_{x,y}E(y)$, with $x$ unused, and
then eliminate this dummy witness.  This leaves exactly $Q_yE(y)$.

For QC first note the fixed theorem
\[
 R:=D_{r,w}(0=0).
\]
RDI at witness $0$ proves $R$: its progressive premise concludes the
reflexive equation $0=0$.  Open $R$ at a fresh variable $b$.  In that
context introduce $D_{x,y}A(y)$ at witness $b$.  This is case mode even
when $A$ is an equation, and the required premise is precisely the given
premise, injectively renamed and weakened.  Eliminate the dummy witness
$x$ to obtain $Q_yA(y)$, and then discharge the opening of $R$; neither
$b$ nor the unused assumption $Q_w(0=0)$ occurs in the result.  Thus QC
does not use a forbidden case introduction of an equation at $0$, even
when the original numerical context is empty.
\end{proof}

Consequently, ordinary induction for a literal equation follows by QP
and closure elimination.  All arguments in the separate predecessor
construction can therefore be performed inside $\mathcal R$ once their
equational derivations have been established.

\subsection{Full formula transport without term substitution on proofs}

\begin{lemma}[Leibniz transport]
For every native formula template $H(v)$, the rule
\[
 \Gamma\vdash s=t,\quad\Gamma\vdash H(s)
 \quad\Longrightarrow\quad\Gamma\vdash H(t)
 \tag{L}
\]
is admissible in $\mathcal R$.
\end{lemma}
\begin{proof}
Induct on the number of formula binders in $H$.  State the induction
hypothesis uniformly for every ambient context and every assignment of
terms to the other free variables.  This uniform statement is essential:
we apply the induction hypothesis to newly formed instances, rather than
substitute terms into previously constructed derivations.

For $H(v)=(a(v)=c(v))$, congruence transports both sides and transitivity
joins the resulting equations.

Suppose $H(v)=Q_yA(v,y)$.  Apply QC to the target closure.  Under its
temporary assumption $A(t,0)$, the reverse equation $t=s$ and the
induction hypothesis give $A(s,0)$.  Eliminate the available closure
$Q_yA(s,y)$ at $Sz$ to obtain $A(s,Sz)$.  The forward equation $s=t$ and
the induction hypothesis give $A(t,Sz)$.  This proves the QC premise.
This explicitly handles the contravariant occurrence of the zero row
in the ordinary meaning of $Q$.

Suppose $H(v)=D_{x,y}B(v,x,y)$.  Open the given $H(s)$ at a fresh
variable $b$, obtaining $Q_yB(s,b,y)$.  Introduce the target $H(t)$ at
the permitted variable witness $b$.  Its case premise assumes
$B(t,b,0)$.  Transport that assumption backwards to $B(s,b,0)$,
eliminate the opened closure at $Sz$, and transport the result forwards
to $B(t,b,Sz)$.  All three formulas have fewer binders than $H$.
RDI at $b$ now proves $H(t)$, and the outer elimination discharges $b$.

The equality premise and other existing derivations are only weakened
or injectively renamed when entering a binder.  No application of an
unproved term-substitution principle occurs.
\end{proof}

\subsection{Naming every term with a fresh variable}

Assume the separately proved equational predecessor lemma: there is a
raw iterator term $p(v)$ such that
\[
 \mathcal R\vdash p(Sv)=v.
 \tag{P}
\]
Then successor injectivity is derived by congruence with $p$ followed by
the two instances of (P).  This step must be justified before the
following argument; importing predecessor from the full iPRA
bi-interpretation would be circular.

Put
\[
 A_v(x,y):=\bigl(C(y,v,C(v,0,Sx))=v\bigr),\qquad
 T(v):=D_{x,y}A_v(x,y).
\]
The defining rows have the derived equations
\[
 A_v(x,0)\;\longleftrightarrow\;v=v,
 \qquad
 A_v(x,Sw)\;\longleftrightarrow\;C(v,0,Sx)=v.
\]
Here and below a displayed equivalence between equations abbreviates
transport along the indicated computation equations.

\begin{lemma}[The naming packet]
For every term $t$, $\mathcal R$ proves $T(t)$.  Moreover, the following
fresh-variable naming rule is admissible:
\[
 \frac{\Gamma,b=t\vdash F}{\Gamma\vdash F},
 \qquad b\notin\mathrm{FV}(\Gamma,t,F).
 \tag{Name}
\]
\end{lemma}
\begin{proof}
First prove $T(0)$ by RDI at witness $0$.  Its matrix is an equation, so
the premise is progressive.  Its successor row reduces by the two
selector computations to $0=0$, independently of the induction
hypothesis.

Next, with $z$ a parameter, prove $T(Sz)$ by RDI at the permitted
variable witness $z$.  Its successor row reduces to
$C(Sz,0,Sz)=Sz$, and then to $Sz=Sz$, independently of its case
assumption.  QC applied to $v\mapsto T(v)$ now gives $Q_vT(v)$.
Closure elimination at the arbitrary term $t$, using $T(0)$, gives
$T(t)$.  This last step is an application of QE, not substitution into
the derivation of $T(v)$.

To prove Name, use the theorem $T(St)$ and open its witness at a fresh
$b$.  The resulting assumption is $Q_yA_{St}(b,y)$.  Its zero row is
a computation theorem, so specialization at $1$ gives
\[
 C(1,St,C(St,0,Sb))=St.
\]
The left side computes first to $C(St,0,Sb)$ and then to $Sb$.
Hence $Sb=St$, and (P) gives $b=t$.  Cut this equation into the given
derivation of $F$, and discharge the opened witness.  The conclusion
does not contain $b$.
\end{proof}

\subsection{Recovery of every omitted introduction}

\begin{theorem}[Restriction of witnesses to zero and variables]
Subject to the equational predecessor lemma (P), $\mathcal R$ and the
native scalar kernel derive exactly the same sequents.  In particular,
the combined witness-mode introduction may be restricted to witnesses
that are either $0$ or variables without changing the theory.
\end{theorem}
\begin{proof}
It suffices to recover the native case introduction at an arbitrary term
$t$; native equational progress was recovered by QP.  Suppose
\[
 \Gamma,B(t,0)\vdash B(t,Sz)
 \tag{*}
\]
with the usual freshness conditions.  Name $t$ by a fresh variable $b$.
Inside the context $\Gamma,b=t$, introduce $D_{x,y}B(x,y)$ at the
allowed witness $b$.  Under its temporary assumption $B(b,0)$, full
formula transport gives $B(t,0)$.  Cut this into a weakened and
injectively renamed copy of (*) to obtain $B(t,Sz)$.  Full formula
transport along $t=b$ gives $B(b,Sz)$.  RDI at $b$ proves the desired
$D$ formula.  Name then removes $b=t$ and the fresh variable $b$.

This argument also recovers the omitted case introduction at literal
$0$ for an equation.  Conversely, every rule of $\mathcal R$ is
derived in the native kernel: case mode is its case introduction,
and progressive mode is obtained from its equation-progress closure
followed by case introduction.  Induction on derivations in either
direction completes the sequent equivalence.
\end{proof}

\begin{corollary}[Bi-interpretability and consistency strength]
The ordinary realization of $\mathcal R$ is bi-interpretable with full
first-order iPRA by the same interpretations as the native scalar
kernel.  Both number domains remain whole, equality remains actual
equality, and both domain composites are the identity graph.  In
particular the theories are equiconsistent.
\end{corollary}
\begin{proof}
The syntax and its ordinary reading have not changed, and the theorem
gives equality of all contextual derivability relations, not merely
equality of closed or arithmetically restricted consequences.  Therefore
the already established native bi-interpretation transfers unchanged.
\end{proof}

\paragraph{Relation to the pinned baseline.}
The README at the pinned commit
\texttt{9f9545f5ffd9},
under ``Further questions'', item 2, conjectures the sufficiency of
case-introduction witnesses $0$ and variables, and correctly identifies
the danger of silently using substitution into derivations of the
restricted calculus.  The proof here resolves that particular logical
gap through QC, the explicit packet $T$, and the independently
equational predecessor lemma.  It does not assert any absolute
minimality of the resulting checker, and its mathematical completeness
claim is separate from any program-size measurement.

\paragraph{Auditable source interfaces.}
The pinned baseline's \texttt{NativeEq.lean} establishes the ordinary
one-term equality interface and full formula transport;
\texttt{NativeGate.lean} and \texttt{NativeLogic.lean} establish the
standard logical encodings; and \texttt{NativeBiInterpretation.lean}
states the native/iPRA comparison.  The argument above reproves the
transport and closure steps in the restricted calculus instead of
importing them with their unrestricted proof-substitution lemmas.
The predecessor construction must likewise be given an independent
equational proof; the statement (P) is displayed to make this
dependency explicit.
